\section{Discussion}
\label{sec:discussion}

\subsection{Where CC-TS pays its conservatism}

The conservative-gate guarantee (\cref{thm:conservative}) is not free. The
policy forgoes savings on cells where the conformal lower bound
$\widehat{q}_{\mathrm{lo}}$ has not yet reached the tolerance threshold
even though the point estimate suggests a downgrade is safe. In our
deployment, this bites hardest on novel agent types --- agents added by
consumer repos that have no offline prior --- and on cells where the eval
suite is silent. The consequence is a measurable opportunity cost
(\cref{fig:ablation-gate}) that an operator might rationalise as
``leaving money on the table.'' The design choice is to take that hit
rather than risk silent regression.

\subsection{Why static frontmatter persists as a safety hatch}

CC-TS is deployed alongside, not in place of, the static-frontmatter
default. Each agent's declared tier is the policy floor: the advisor may
recommend a downgrade, but the operator can lock the cell to the static
default via a single configuration flag. This is more than belt-and-braces
caution. Two regimes specifically motivate the safety hatch:

\begin{enumerate}[leftmargin=*]
  \item \textbf{Incident response.} During an active incident, exploration
        is exactly the wrong policy. The advisor must yield to the
        operator's risk tolerance, not the calibrated default. We expose a
        \texttt{force\_baseline=true} flag that short-circuits all learning
        and plays the static recommendation.
  \item \textbf{Adversarial drift.} If the upstream model provider
        deprecates or silently changes a tier, posterior estimates become
        stale before the advisor's change-point detection can react. The
        static floor caps the damage.
\end{enumerate}

\subsection{Off-policy evaluation gap}

The cost results in \cref{sec:eval} use doubly-robust off-policy evaluation
\cite{dudik2014doubly} to compare CC-TS against baselines that were never
actually run on the production trace. This is the standard approach in
contextual-bandit deployments \cite{agarwal2016making}, but it has a
known limitation: when the dispatch policy and the reward distribution are
both non-stationary --- which is exactly our setting, since both the
advisor's posteriors and Anthropic's underlying models are evolving --- DR
estimators bias toward whichever regime dominates the logged data. We
report the bootstrap CIs honestly, but operators reading the paper should
understand that the cost numbers are point estimates with non-trivial
distributional assumptions baked in.

\subsection{When this won't work}

Three regimes break CC-TS as currently described:

\begin{itemize}[leftmargin=*]
  \item \textbf{Continuous quality signal.} If the eval suite emits a
        continuous score in $[0, 1]$ rather than a binary pass/fail
        indicator, the Beta-Bernoulli posterior is the wrong family. A
        Gaussian (or Beta-distributed continuous-quality) posterior is the
        natural extension; this is open question OQ-A6 in
        \cite{blackrim_moa1c}.
  \item \textbf{Sequential structure.} CC-TS treats each dispatch as a
        one-shot decision. If a sequence of dispatches has joint downstream
        consequences --- e.g., an architect's ADR cascades into seven
        builder dispatches --- the per-cell posterior under-counts the
        cascade cost. The critical-path-aware extension in
        \cite{blackrim_moa10} addresses this; the formal treatment is
        contextual MDP, not bandit.
  \item \textbf{Multi-tenant federation.} A consumer repo running CC-TS in
        isolation sees only its own telemetry. A federated CC-TS that
        shares posteriors across tenants would dramatically improve
        thin-cell convergence, but introduces privacy and trust questions
        out of scope for this paper.
\end{itemize}

\subsection{Lessons from deployment}

Three deployment lessons are worth flagging explicitly because they
generalise beyond CC-TS to any production bandit deployment.

\paragraph{Quality observation is the hard part.}
The bandit literature is dense; the production-quality-signal literature is
sparse. Our three-channel observation mechanism (\cref{sec:system}) was
the most carefully-designed part of the system, and the most fragile. The
disk-flush race in transcript persistence \cite{blackrim_hyqd} alone cost
us a week of debugging. We commend the engineering effort to anyone
attempting a similar deployment: the bandit policy is the easy part; the
reward signal is the hard part.

\paragraph{Audit surface beats accuracy.}
Operators trust an advisor whose recommendations they can read and check.
The structured \texttt{Reasons} object on every recommendation
(\cref{sec:layer3}) was responsible for more operator-trust uplift than
any improvement to the underlying algorithm. A black-box advisor with
$10\%$ better policy accuracy would have been turned off in week one.

\paragraph{Conservative bias is the design contract.}
The asymmetric-loss multiplier $M[\text{Critical}] = \infty$
(\cref{tab:loss-matrix}) collapses to a hard rule: critical-path cells
never enter the exploration set. This is not a hyperparameter; it is the
design contract under which operators allow the advisor to drive
production routing. Loosening it would surrender the trust that the
deployment is built on.
